%Este es documento para el libro del proyecto
\documentclass[11pt,openany,oneside,letterpaper]{book}
\usepackage[spanish]{babel}
\usepackage[utf8]{inputenc}
\usepackage{url}

%%Dimensiones de p�gina
\usepackage[hmargin={3.0cm,2.5cm},vmargin={3.0cm,2.5cm}]{geometry}
\usepackage[spanish]{layout}
\marginparwidth=2.0cm
%\layout %Para mirar como esta configurada la p�gina

%Sangria y espacio entre p�rrafos
%\parindent=0pt
\parskip=10pt

%Paquetes empleados
\usepackage{array,amsmath,amssymb,amsfonts,pifont,verbatim}
\usepackage{dsfont}

%Paquetes para el formato del libro
\usepackage{fancyhdr}
\usepackage[]{fncychap}

\usepackage{pstricks,pst-all}

\usepackage[dvips]{graphicx}

%Paquete para insertar mejor los objetos flotantes
\usepackage{float}



\newenvironment{abstract}{%
    \chapter*{Resumen}%
    \addcontentsline{toc}{chapter}{Resumen}%
}{%
}
%Numeracion de secciones bajo capitulos
\renewcommand{\thesection}{\thechapter.\arabic{section}}


%Redefinition de t�tulos y r�tulos
\renewcommand{\listtablename}{\'Indice de tablas}
\renewcommand{\tablename}{Tabla}


\begin{document}
\frontmatter

\tableofcontents



\mainmatter

%Parte de la formulacion del proyecto
\part{Descripcion del proyecto}

\chapter{Descripción del proyecto}
\section{Resumen}
La transcripción automática de música (Automatic Music Transcription, AMT) es un área de investigación dentro del campo de Music Information Retrieval (MIR) cuyo objetivo consiste en transformar señales de audio musical en representaciones simbólicas como notas, acordes o partituras. Este problema resulta especialmente complejo en instrumentos polifónicos como la guitarra acústica, donde múltiples notas se producen simultáneamente generando espectros armónicos superpuestos.

En este trabajo se presenta un enfoque para la detección e identificación automática de acordes de guitarra acústica a partir de grabaciones de audio. La metodología propuesta se basa en el procesamiento espectral de la señal y en el uso de modelos de aprendizaje automático para reconocer patrones armónicos asociados a distintos acordes musicales.

Este estudio busca contribuir al desarrollo de herramientas computacionales que faciliten el análisis automático de interpretaciones musicales y la generación de representaciones simbólicas a partir de señales de audio~\cite{benetos2019automatic}.
\section{Justificación}

La transcripción automática de la música presenta un desafío al momento de procesar las señales digitales, de igual forma constituye un problema que toma gran relevancia al momento de recuperar información musical (MIR), esta tiene un gran potencial para transformar la manera en que se analiza, se interpreta y se accede al contenido musical~\cite{jamshidi2024machine}. A pesar de los avances recientes en los ámbitos de machine learning, lo sistemas actuales no tienen la capacidad máxima de precisión y flexibilidad al comparar con un músico profesional, ya que estos tienen la capacidad de entender información polifónica mucho más amplia, en estos escenarios polifónicos complejos se convergen múltiples voces, sonidos y narrativas sonoras que son complejas para una maquina de entender de manera simultanea~\cite{benetos2019automatic}.

El desarrollo de aplicaciones que permitan la extracción de información como acordes, estructuras musicales, datos de frecuencia y demás, puede llegar a ser una herramienta clave para un aprendizaje musical, que puede acompañar al interesado y así darle un respectivo uso educativo. La base para el desarrollo de una aplicación parte de una transcripción musical, donde permite la conversión de información auditiva en representaciones simbólicas, esto constituye la base de un estudio, una interpretación y una aplicación en respectivos contextos~\cite{jamshidi2024machine}. En este sentido, automatizar este proceso contribuye en una reducción del tiempo y de dificultad a interpretadores cuya tarea sea transcribir de forma manual y presentarle más información que no sea solo en formato de partitura.


\section{Introducción}
\subsection{Descripción del problema}

En los últimos años, la inteligencia artificial ha avanzado hasta resolver tareas que hace una década parecían prohibitivas. La detección de objetos en tiempo real, el análisis de imágenes médicas y la traducción automática son problemas que hoy tienen soluciones competitivas, a veces superiores al desempeño humano en condiciones controladas. Sin embargo, no todo el territorio está cubierto: la música resiste, y lo hace por razones que no se resuelven añadiendo más datos o más capas, como la percepción, el contexto y las decisiones interpretativas, que no caben fácilmente en una función de pérdida.

La transcripción musical automática (AMT) sigue siendo un reto, ya que va más allá del procesamiento y análisis de audio para su posterior transcripción. En el caso de los instrumentos polifónicos, se debe tener en cuenta la superposición espectral que se produce al formar acordes, lo que conlleva consideraciones más detalladas, como el material de las cuerdas. En la guitarra acústica con cuerdas de nailon, por ejemplo, se presenta una menor inarmonicidad y un \textit{decay} más pronunciado que en las cuerdas metálicas, lo que altera la distribución de armónicos en el espectrograma.

Además de las condiciones físicas del instrumento, se deben tener en cuenta las distintas técnicas de ejecución. Por ejemplo, un Am7 puede sonar perceptiblemente distinto según se toque con dedos o púa, en distintas posiciones del mástil, con vibrato, \textit{muting} parcial o cambios de dinámica. De igual forma, conceptos de la teoría musical, como la enarmonía, introducen otra capa de complejidad. Do\# y Re$\flat$ comparten frecuencia fundamental, pero representan acordes distintos según el contexto armónico, una distinción que el espectrograma por sí solo no puede resolver y que requiere contexto adicional o intervención humana para mejorar su precisión.

\chapter{Objetivos}
\input{src/4. Objetivos}


\part{Marco teorico}

\chapter{Marco de Referencia}
\section{Estado del arte}

La transcripción automática de música (Automatic Music Transcription, AMT) es una de las áreas más importantes dentro del campo de \textit{Music Information Retrieval} (MIR). Para simplificar, el objetivo de la transformación es convertir lo que escuchamos en una canción en algo que se pueda escribir, leer e interpretar, como lo son las notas musicales o una partitura en específico \cite{benetos2018amt}. Parece algo sencillo de entender, pero en realidad es un problema bastante complejo y en donde no se ha tocado fondo; esto se debe a que una señal de audio puede contener múltiples sonidos ocurriendo al mismo tiempo, lo que al querer identificar correctamente cada nota se torna muy problemático: cuándo empieza, cuándo termina y de dónde proviene.

Indagando en los principios de la transcripción automática de música, la mayoría de los enfoques que se toman para resolver este problema se basan en técnicas clásicas de procesamiento de señales. Estas técnicas utilizaban transformadas espectrales y modelos probabilísticos para intentar identificar eventos musicales a partir del espectrograma \cite{klapuri2006transcription}. En sí, este tipo de métodos funcionaba relativamente bien, pero únicamente cuando se trataba de una sola nota a la vez; esto se conoce como señales monofónicas. Sin embargo, esto presentaba muchas dificultades cuando varias notas sonaban simultáneamente, como ocurre en los acordes, debido a la superposición de sus componentes frecuenciales en cada caso.

Con respecto al problema de la superposición, con el tiempo los investigadores comenzaron a explorar representaciones más adecuadas para la música. Una de las más importantes es la \textit{Constant-Q Transform} (CQT), que permite representar las frecuencias de una forma más cercana a cómo las percibe el oído humano, haciendo de esta una herramienta clave para el uso y el trabajo de los acordes \cite{sigtia2020representation}. Gracias a esto, los sistemas más actuales y usados pueden reconocer mejor patrones armónicos, notas y frecuencias mucho más avanzadas, lo cual es fundamental para el tema de identificar acordes y estructuras musicales más complejas.

Posteriormente al CQT, apareció un fenómeno mundial que definiría el futuro de la era tecnológica: el aprendizaje automático. La incorporación del ML marcó un punto de inflexión en el estudio futuro en este campo. En lugar de depender únicamente de reglas matemáticas, los modelos comenzaron a aprender directamente a partir de los datos que eran retroalimentados constantemente. Como ejemplo, se ha visto que las redes neuronales convolucionales pueden identificar patrones característicos en las señales de audio que corresponden a eventos musicales específicos, lo que facilita en gran medida el trabajo del tema \cite{cnn_amt2025}. De una manera muy similar, otros trabajos que se han hecho han mostrado mejoras en la precisión de la transcripción al analizar características acústicas, esto llevado a cabo mediante algoritmos de aprendizaje automático que monitorean y hacen uso de datos reales \cite{jaciii2025}.

Ya para tiempos más recientes, los avances en aprendizaje profundo han permitido desarrollar modelos más sofisticados capaces de abordar la transcripción musical de forma más íntegra con la capacidad humana de interpretación. Un ejemplo de esto es el enfoque \textit{end-to-end}, donde consiste en que el sistema recibe directamente el audio y produce la transcripción sin necesidad de etapas intermedias complejas que alarguen y complejicen más el caso \cite{hawthorne2019polyphonic}. Este tipo de modelos es especialmente útil en instrumentos como la guitarra, como es el caso de estudio, donde los acordes están formados por varias notas que se superponen en el espectro de frecuencias y son de necesario abordaje y estudio.

Además de lo anterior, algunos estudios han propuesto mejorar estos sistemas incorporando información adicional que nutra mejor los datos. Por ejemplo, el uso de estadísticas musicales que permitan hacer transcripciones más coherentes desde el punto de vista estructural y analítico de un sistema avanzado \cite{nonlocal2021}. De igual forma, se han desarrollado soluciones específicas para distintos tipos de notación o instrumentos de gran complejidad, como tablaturas de órgano o en instrumentos tradicionales, lo que de por sí demuestra la diversidad que hay y los enfoques dentro del área que pueden ser abordados \cite{organ2022, qanun2024}.

En otra línea de trabajo que se puede abordar y es de suma importancia, está la adaptación de los modelos a contextos musicales específicos. Técnicas como el \textit{fine-tuning} permiten ajustar modelos previamente entrenados para que reconozcan mejor ciertos estilos o características particulares de la música \cite{musicnet2022finetuning}. Como también ocurre, se ha explorado el uso de información contextual para mejorar la interpretación de las señales según sea el contexto, lo cual es clave en un dominio donde el significado musical depende en gran medida de un entorno en el que se desarrolla \cite{ins2021}.

Los datos y revisiones recientes coinciden en que los mejores resultados se logran combinando buenas representaciones de una señal, modelos de aprendizaje profundo y grandes conjuntos de datos etiquetados, siempre y cuando se manejen los datos con mucho detalle y rigurosidad \cite{survey2024}. En esta línea de trabajo, varios proyectos actuales utilizan directamente espectrogramas como entrada a modelos de redes neuronales, logrando avances significativos en la detección de notas, acordes y progresiones armónicas en mucho menos tiempo y con una efectividad enorme \cite{cqt2025, melody2025}.

Ya para algo más general, la evolución de este entorno, del área en sí, muestra un paso claro desde métodos tradicionales hacia modelos basados en aprendizaje profundo. Este representa un cambio que ha permitido mejorar considerablemente el rendimiento de los sistemas de transcripción automática; a pesar de todos los avances, todavía existen desafíos importantes, especialmente en escenarios polifónicos y en la generalización a distintos instrumentos que pueden representar una dificultad mucho mayor a uno corriente y que, desde luego, son aspectos que se quiere llegar a dominar.

Finalmente, en la búsqueda de más información se encuentran otros estudios que han abordado problemas relacionados, como la alineación entre audio y partituras, donde se busca, mediante la utilización de redes neuronales recurrentes, la mejora significativa de la precisión de la transcripción de la música \cite{audioalign2018}. También se han desarrollado sistemas automáticos de generación de partituras y modelos híbridos que combinan análisis espectral con aprendizaje automático y que ya dan un paso más adelante en la comprensión total \cite{jcss2021, jme2024}. Viéndolo desde una perspectiva más amplia, algunos trabajos destacan la importancia de integrar conocimientos de teoría musical con enfoques computacionales y son el eje principal en el desarrollo de sistemas que comprendan la música \cite{digitalhumanities2022}; y, por otra parte, están las revisiones del área que identifican como principales retos la complejidad de la polifonía y la adaptación de los modelos a diferentes contextos musicales, ampliando por mucho la comprensión del tema \cite{jiis2014}.
\input{src/8. Marco teorico}
\part{Desarrollo}
\chapter{Diseño e implementación}
\input{src/5. Alcance}
\section{Limitaciones}
A pesar de los avances en el campo de la transcripción automática de música, existen varias limitaciones que deben ser consideradas al desarrollar y evaluar sistemas de detección de acordes. En primer lugar, la calidad y diversidad del conjunto de datos utilizado para entrenar los modelos pueden afectar significativamente su rendimiento. Si el conjunto de datos no incluye una amplia variedad de estilos musicales, técnicas de ejecución y condiciones de grabación, el modelo puede tener dificultades para generalizar a situaciones del mundo real.Además, la presencia de ruido de fondo, variaciones en la afinación y técnicas de ejecución no convencionales pueden degradar el rendimiento del sistema, lo que limita su aplicabilidad en entornos no controlados~\cite{benetos2019automatic}. 

Otra limitación importante es la capacidad de los modelos para manejar la polifonía compleja de la guitarra acústica, donde múltiples notas se superponen y generan espectros armónicos complejos. Finalmente, aunque se pueden lograr altos niveles de precisión en la detección de acordes, es importante reconocer que ningún sistema de transcripción automática puede igualar completamente la capacidad de un músico profesional para interpretar y comprender la música, lo que implica que siempre habrá un margen de error y limitaciones en la interpretación de los resultados~\cite{jamshidi2024machine}.
\section{Plan de trabajo}

Para el desarrollo de este trabajo de investigación se propone el siguiente plan de trabajo, dividido en etapas con sus respectivas actividades y cronograma, el cual se desarrollará durante un período de dos semestres académicos y se detallará en su respectiva sección más adelante:

\subsection{Revisión documental}

Durante esta etapa, se realizará la búsqueda y recopilación de datasets de audio de guitarra acústica, con la mayor cantidad de información adicional posible como notaciones o archivos MIDI para mejorar la calidad de las etiquetas y garantizar la mejor calidad de los datos posible.

\subsection{Implementación del modelo}
En esta fase, se llevará a cabo la implementación de los modelos de aprendizaje automático seleccionados para la detección de acordes, empleando prácticas de codificación eficientes y asegurando la reproducibilidad de los resultados. Se documentará el proceso de implementación y se prepararán los scripts necesarios para el entrenamiento y evaluación de los modelos.

\subsection{Preprocesamiento de datos}

En esta fase, se llevará a cabo la limpieza y normalización de los datos del o los datasets elegidos para el proyecto, así como la extracción de características relevantes para la detección de acordes. Se evaluarán diferentes técnicas de preprocesamiento para mejorar la calidad de las representaciones de audio.

\subsection{Entrenamiento del modelo}

En esta etapa, se implementarán y entrenarán diferentes modelos de aprendizaje automático para la detección de acordes, utilizando las características extraídas en la fase anterior. Se optimizarán los hiperparámetros de los modelos para mejorar su rendimiento.

\subsection{Comparación de modelos}

Una vez entrenados los modelos, s.e realizará una comparación exhaustiva de su desempeño utilizando métricas de evaluación como exact‑match accuracy, top‑3 accuracy y macro‑F1. Se analizarán los resultados para determinar cuál modelo ofrece el mejor rendimiento en la tarea de detección de acordes.

\subsection{Pruebas}

Finalmente, se evaluará la robustez del mejor modelo frente a variaciones de afinación, técnicas de ejecución y presencia de ruido de fondo, comparando los resultados obtenidos con una revisión manual bajo los criterios de teoría musical. Se documentarán los resultados obtenidos y se discutirán las limitaciones del sistema propuesto, así como posibles mejoras y aplicaciones futuras en el ámbito de la transcripción musical automática y el análisis de audio.

\newpage

\section{Cronograma}

El cronograma propuesto para el desarrollo de la investigación está distribuido en un período de ocho (8) meses a partir de la aprobación del anteproyecto. El Cuadro \ref{tab:cronograma} ilustra la distribución temporal de las fases del plan de trabajo.

\begin{table}[H]
\centering
\caption{Cronograma de actividades del proyecto}
\label{tab:cronograma}
\begin{tabular}{|l|c|c|c|c|c|c|c|c|}
\hline
\textbf{Actividad} & \textbf{M1} & \textbf{M2} & \textbf{M3} & \textbf{M4} & \textbf{M5} & \textbf{M6} & \textbf{M7} & \textbf{M8} \\
\hline
Revisión documental & $\bullet$ & $\bullet$ & & & & & & \\
\hline
Preprocesamiento de datos & & $\bullet$ & $\bullet$ & & & & & \\
\hline
Implementación del modelo & & & & $\bullet$ & $\bullet$ & & & \\
\hline
Entrenamiento del modelo & & & & & $\bullet$ & $\bullet$ & $\bullet$ & \\
\hline
Comparación de modelos & & & & & & & $\bullet$ & \\
\hline
Pruebas & & & & & & & & $\bullet$ \\
\hline
\end{tabular}
\end{table}

Al finalizar cada una de estas actividades, se espera cumplir con los siguientes entregables principales:

\begin{itemize}
\item \textbf{Mes 1 - 2 | Revisión documental:} Datasets seleccionados y documentados, junto con la revisión bibliográfica consolidada.
\item \textbf{Mes 2 - 3 | Preprocesamiento de datos:} Pipeline de código funcional con la extracción de características espectrales validadas.
\item \textbf{Mes 4 - 5 | Implementación del modelo:} Arquitecturas base implementadas y listas para su entrenamiento.
\item \textbf{Mes 5 - 7 | Entrenamiento del modelo:} Modelos entrenados con hiperparámetros optimizados y resultados preliminares registrados.
\item \textbf{Mes 7 | Comparación de modelos:} Reporte de desempeño comparativo con las métricas establecidas (exact-match accuracy, top-3 accuracy, macro-F1).
\item \textbf{Mes 8 | Pruebas:} Validación de robustez finalizada y entrega del documento de tesis.
\end{itemize}
\section{Desarrollo}
Para el desarrollo de este trabajo se utilizará principalmente herramientas de Python, como lo son las librerías \textit{librosa} y \textit{pytorch}, las cuales ofrecen una amplia gama de funcionalidades para el procesamiento de audio y la extracción de características musicales. Estas librerías permiten realizar tareas como la lectura de archivos de audio, la generación de espectrogramas, la aplicación de la Transformada de Fourier y la extracción de características cromáticas, entre otras.
Como avance entregable se presentará el entrenamiento y procesamiento de datos con la implementación de un modelo de aprendizaje CNN 2‑D, el cual se entrenará utilizando un dataset de audio de guitarra acústica de \textit{Guitarset} con etiquetas de acordes. Donde se evaluará el desempeño del modelo utilizando métricas como exact‑match accuracy, top‑3 accuracy y macro‑F1, con el objetivo de alcanzar una precisión exacta mínima del 85 $\%$ en la identificación de acordes. Además, se documentará el proceso de desarrollo y se prepararán los scripts necesarios para la reproducción de los resultados obtenidos, así como también una pequeña interfaz de usuario para facilitar su uso.


\backmatter


\bibliographystyle{IEEEtran}
\bibliography{references}
\end{document}